% Generated by experiments/scripts/export_latex.py -- do not edit by hand.
% Source: experiments/data/2-processed-tables-and-evaluations/revised-usd/holdout-w0/eval_summary.json
\begin{table}[htbp]
\centering
\caption[스펜더 코호트의 시간 홀드아웃 W0.]{스펜더 코호트(계약 $n=17{,}050$개)에서 해 본 시간 홀드아웃 W0이에요. 2026년 3월 31일에 동결한 점수와, 2026년 4월 1일부터 2026년 6월 30일까지 그 계약에서 센 두 라벨 사이의 켄달의 $\tau$를 보여 줘요. 두 라벨은 새로 생긴 주인--스펜더 허락 쌍의 수와, 그 계약에 처음으로 송금을 보낸 주소의 수예요. 코호트에는 동결일에 마지막 허락이 0보다 큰 스펜더 17{,}167개 가운데 계약 17{,}050개를 남겼어요. 코드가 없는 EOA 104개와 EIP-7702 위임이 있는 EOA 13개는 뺐어요. 계약 2{,}070개는 새 쌍을 하나 이상 받았고, 1{,}695개는 새 송금자를 하나 이상 받았어요. 기준선 행은 각 한 층 페이지랭크가 매끄럽게 다듬는 $t_1$의 원시 차수예요. C-PR ($\lambda=1$)과 C-PR ($\lambda=0$)은 결합 연산자의 허락 층과 송금 층을 각각 한 층만 걷는 것이에요. 신뢰구간(CI)은 계약을 다시 뽑아 구한 95\% 백분위 부트스트랩 구간이에요(짝지은 다시 뽑기 400번, 시드 42).}
\label{tab:holdout-spenders}
\fitwidth{%
\begin{tabular}{lcc}
\toprule
$t_1$의 점수 & 새 허락 쌍 ($\tau$ [95\% 신뢰구간]) & 새 송금자 ($\tau$ [95\% 신뢰구간]) \\
\midrule
$t_1$의 받은 허락 수 (기준선) & 0.250 [0.237, 0.262] & 0.180 [0.169, 0.192] \\
$t_1$의 송금 입차수 (기준선) & 0.135 [0.121, 0.150] & 0.303 [0.292, 0.315] \\
\midrule
EndorseRank & 0.231 [0.220, 0.242] & 0.200 [0.189, 0.210] \\
AWP & 0.196 [0.181, 0.210] & 0.346 [0.336, 0.356] \\
\midrule
C-PR ($\lambda=1$) & 0.223 [0.213, 0.233] & 0.188 [0.177, 0.199] \\
C-PR ($\lambda=0$) & 0.160 [0.144, 0.174] & 0.296 [0.285, 0.307] \\
C-PR ($\lambda=0.25$) & 0.253 [0.243, 0.264] & 0.269 [0.257, 0.280] \\
C-PR ($\lambda=0.5$) & 0.252 [0.241, 0.262] & 0.262 [0.250, 0.273] \\
C-PR ($\lambda=0.75$) & 0.249 [0.238, 0.260] & 0.253 [0.242, 0.265] \\
S-PR & 0.172 [0.157, 0.185] & 0.295 [0.284, 0.305] \\
\bottomrule
\end{tabular}
}
\end{table}
